\documentclass[11pt,spanish]{article}

% =========================================================
% PLANTILLA BASE PARA TAREAS ACADÉMICAS / INFORMES TÉCNICOS
% =========================================================

% --- Idioma y codificación para Tectonic / XeLaTeX ---
\usepackage[spanish,es-tabla,es-noshorthands]{babel}
\usepackage{fontspec}

% --- Márgenes / papel ---
\usepackage{geometry}
\geometry{a4paper, margin=2.7cm}

% --- Matemáticas ---
\usepackage{amsmath,amssymb}

% --- Tablas, enlaces y elementos visuales ---
\usepackage{array}
\usepackage{tabularx}
\usepackage{booktabs}
\usepackage{longtable}
\usepackage{graphicx}
\usepackage{float}
\usepackage{xcolor}
\usepackage{tikz}
\usepackage{pgfplots}
\pgfplotsset{compat=1.16}
\usetikzlibrary{positioning,shapes.geometric,arrows.meta,calc}
\usepackage{enumitem}
\usepackage{ragged2e}
\usepackage{caption}
\captionsetup{skip=6pt}
\captionsetup[table]{position=top}
\captionsetup[figure]{position=bottom}
\usepackage{xurl}
\usepackage[hidelinks]{hyperref}

% --- Configuración de párrafo en bloque ---
\makeatletter
\renewcommand\paragraph{\@startsection{paragraph}{4}{\z@}%
  {2.2ex \@plus1ex \@minus.2ex}%
  {0.8ex}%
  {\normalfont\normalsize\bfseries}}
\makeatother

% --- Listados de código con formato elegante ---
\usepackage{listings}

\definecolor{codebg}{rgb}{0.97,0.97,0.98}
\definecolor{codeframe}{rgb}{0.85,0.85,0.88}
\definecolor{commentgreen}{rgb}{0.13,0.54,0.13}
\definecolor{keywordblue}{rgb}{0.0,0.2,0.7}
\definecolor{stringpurple}{rgb}{0.6,0.1,0.4}

% --- Colores para esquemas y diagramas TikZ ---
\definecolor{softblue}{rgb}{0.88, 0.93, 1.0}
\definecolor{softgreen}{rgb}{0.88, 0.96, 0.88}
\definecolor{softorange}{rgb}{1.0, 0.92, 0.82}
\definecolor{softpurple}{rgb}{0.93, 0.88, 0.98}
\definecolor{softteal}{rgb}{0.85, 0.95, 0.95}
\definecolor{softyellow}{rgb}{1.0, 0.95, 0.80}

\lstdefinestyle{customcode}{
    backgroundcolor=\color{codebg},
    basicstyle=\ttfamily\footnotesize,
    breakatwhitespace=false,
    breaklines=true,
    captionpos=b,
    keepspaces=true,
    numbers=none,
    showspaces=false,
    showstringspaces=false,
    showtabs=false,
    tabsize=4,
    frame=single,
    rulecolor=\color{codeframe},
    commentstyle=\color{commentgreen}\itshape,
    keywordstyle=\color{keywordblue}\bfseries,
    stringstyle=\color{stringpurple},
    xleftmargin=0.5em,
    xrightmargin=0.5em
}

\lstdefinestyle{customcodenum}{
    backgroundcolor=\color{codebg},
    basicstyle=\ttfamily\footnotesize,
    breakatwhitespace=false,
    breaklines=true,
    captionpos=b,
    keepspaces=true,
    numbers=left,
    numbersep=6pt,
    numberstyle=\tiny\color{gray},
    showspaces=false,
    showstringspaces=false,
    showtabs=false,
    tabsize=4,
    frame=single,
    rulecolor=\color{codeframe},
    commentstyle=\color{commentgreen}\itshape,
    keywordstyle=\color{keywordblue}\bfseries,
    stringstyle=\color{stringpurple},
    xleftmargin=1.5em,
    xrightmargin=0.5em
}

\lstset{style=customcode}

% --- Encabezado y pie de página: solo número de página en el pie ---
\usepackage{fancyhdr}
\setlength{\headheight}{15pt}
\setlength{\footskip}{25pt}

\pagestyle{fancy}
\fancyhf{}
\renewcommand{\headrulewidth}{0pt}
\renewcommand{\footrulewidth}{0pt}
\fancyfoot[C]{\thepage}

\fancypagestyle{plain}{%
    \fancyhf{}%
    \renewcommand{\headrulewidth}{0pt}%
    \renewcommand{\footrulewidth}{0pt}%
    \fancyfoot[C]{\thepage}%
}

% --- \input de filas generadas dentro de tabular (el \input de LaTeX añade ganchos que rompen \bottomrule) ---
\makeatletter\newcommand{\filas}[1]{\@@input #1 }\makeatother

% --- Interlineado ---
\linespread{1.2}
\sloppy

% =========================================================
% DATOS DEL DOCUMENTO  (Progra02: lo de Persona C)
% =========================================================
\newcommand{\Institucion}{Instituto Tecnológico de Costa Rica}
\newcommand{\Campus}{Escuela de Ingeniería en Computación}
\newcommand{\DocumentoTitulo}{Progra02}
\newcommand{\DocumentoSubtitulo}{Búsqueda y evaluación}
\newcommand{\NombreBuscador}{Encuadre}
\newcommand{\Curso}{IC-8060 Recuperación de Información Textual}
\newcommand{\Profesor}{Aurelio Sanabria}
\newcommand{\FechaDocumento}{\pendiente{fecha de entrega}}
\newcommand{\PeriodoAcademico}{II Semestre}
\newcommand{\AnioDocumento}{2026}

% Lo que falta por llenar se ve en rojo: buscar "PENDIENTE" antes de entregar.
\newcommand{\pendiente}[1]{\textcolor{red}{\textbf{[PENDIENTE: #1]}}}

\newcommand{\AutoresPortada}{%
\begin{minipage}{0.85\textwidth}
\begin{tabular*}{\linewidth}{@{\extracolsep{\fill}}lr@{}}
Santiago Espinoza Rendon & 2024156530 \\
Fabián Araya Ortega & 2019159110 \\
Mauricio González Prendas & 2024143009 \\
José Fabián Zumbado Ruiz & \pendiente{carné} \\
\end{tabular*}
\end{minipage}%
}

\title{\DocumentoTitulo}
\author{Santiago Espinoza Rendon, Fabián Araya Ortega, Mauricio González Prendas, José Fabián Zumbado Ruiz}
\date{\FechaDocumento}

\graphicspath{{../../Progra01/imgs/}{imgs/}}
\usepackage{pgfplotstable}
\pgfplotstableset{col sep=comma}
\pgfplotsset{compat=1.16, width=\linewidth, height=6.5cm, grid=major,
  every axis plot/.append style={thick}}

\newcolumntype{Y}{>{\RaggedRight\arraybackslash}X}
\renewcommand{\arraystretch}{1.22}

\newcommand{\MakeCover}{%
\begin{titlepage}
\thispagestyle{empty}
\centering
{\bfseries \huge \Institucion \par}
\vspace{0.25cm}
{\LARGE \Campus \par}
\vspace{0.15cm}
{\LARGE \Curso \par}
\vspace{0.5cm}
\includegraphics[width=0.22\linewidth]{logo.png}
\vfill
{\huge \bfseries \DocumentoTitulo: \DocumentoSubtitulo \par}
\vspace{0.4cm}
{\Large \textit{\NombreBuscador} \par}
\vfill
{\bfseries \LARGE Profesor: \par}
\vspace{0.2cm}
{\Large \Profesor \par}
\vfill
{\bfseries \LARGE Estudiantes: \par}
\vspace{0.3cm}
{\Large \AutoresPortada \par}
\vfill
{\bfseries\LARGE Fecha: \par}
\vspace{0.2cm}
{\Large \FechaDocumento \par}
\vfill
{\LARGE \PeriodoAcademico \par}
{\LARGE \AnioDocumento \par}
\end{titlepage}%
}

\begin{document}
\hypersetup{pageanchor=false}
\MakeCover
\hypersetup{pageanchor=true}
\setcounter{page}{1}

% =========================================================
\section{Estadísticas de uso}
% =========================================================
Las cifras de esta sección se generan con \texttt{python -m metricas --documento documento}
a partir de los juicios de relevancia guardados por la interfaz de evaluación, de modo que
el documento y la página de métricas siempre muestran los mismos números.
\pendiente{correr el comando con los juicios reales; mientras tanto las tablas y figuras salen vacías}

\subsection{Participación y juicios}

\begin{table}[H]
  \centering
  \begin{tabular}{@{}lr@{}}
    \toprule
    \textbf{Indicador} & \textbf{Valor} \\
    \midrule
    \filas{tablas/uso.tex}
    \bottomrule
  \end{tabular}
  \caption{Participación en las sesiones de evaluación.}\label{tab:uso}
\end{table}

\begin{figure}[H]
  \centering
  \begin{tikzpicture}
    \begin{axis}[ybar, height=5cm, width=0.7\linewidth, bar width=22pt, ymin=0,
      xtick=data, xticklabels from table={datos/juicios_por_evaluador.csv}{evaluador},
      ylabel={Juicios}, nodes near coords, enlarge x limits=0.3,
      enlarge y limits={upper, value=0.2}]
      \addplot[fill=teal!70] table[x=n, y=juicios, col sep=comma]{datos/juicios_por_evaluador.csv};
    \end{axis}
  \end{tikzpicture}
  \caption{Juicios guardados por evaluador.}
\end{figure}

\pendiente{describir quiénes evaluaron (perfil, cuántas sesiones, cómo se coordinaron) y
adjuntar una captura de la interfaz de evaluación}

\subsection{Métricas utilizadas y su justificación}

Con una colección sin juicios previos hay que construirlos con personas, y las métricas deben
elegirse según lo que esos juicios permiten medir~\cite{sanderson2010}. Se usan cuatro,
definidas para una consulta con ranking $d_1,\dots,d_n$, relevancia binaria $r_i$ y
calificación $g_i$.

\begin{table}[H]
  \centering
  \begin{tabularx}{\linewidth}{@{}l>{\RaggedRight}p{5.2cm}Y@{}}
    \toprule
    \textbf{Métrica} & \textbf{Definición} & \textbf{Por qué se eligió} \\
    \midrule
    P@$k$ &
      $\displaystyle \frac{1}{k}\sum_{i=1}^{k} r_i$ &
      Mide lo que ve el usuario en la primera página de resultados; es simple de interpretar
      y no necesita conocer todos los relevantes~\cite{manning2008}. \\
    MAP &
      $\displaystyle \frac{1}{|Q|}\sum_{q}\frac{1}{|R_q|}\sum_{i} P@i\cdot r_i$ &
      Resume la calidad del orden en toda la lista en un solo número y es una medida
      de uso muy extendido para comparar sistemas~\cite{manning2008}. \\
    nDCG@$k$ &
      $\displaystyle \frac{\sum_{i=1}^{k} g_i/\log_2(i+1)}{\text{DCG ideal}}$ &
      Aprovecha la relevancia graduada (0, 1, 2) y descuenta lo relevante que aparece
      abajo~\cite{jarvelin2002}. \\
    MRR &
      $\displaystyle \frac{1}{|Q|}\sum_{q}\frac{1}{\text{pos. del primer relevante}}$ &
      Refleja consultas donde basta un buen resultado~\cite{voorhees1999}. \\
    \bottomrule
  \end{tabularx}
  \caption{Métricas de calidad del ranking.}\label{tab:def-metricas}
\end{table}

El acuerdo entre evaluadores se mide con el coeficiente $\kappa$ de
Cohen~\cite{cohen1960}, $\kappa=(p_o-p_e)/(1-p_e)$, y se lee con la escala (orientativa) de Landis y
Koch~\cite{landis1977}. Cuando dos o más personas juzgan el mismo documento, la calificación
se promedia para nDCG y se considera relevante si al menos la mitad lo marcó $\ge 1$.
Que los evaluadores discrepen no invalida la comparación entre sistemas~\cite{voorhees2000}.

\pendiente{justificar con criterio propio por qué estas cuatro y no otras, a la luz de la
necesidad de información del grupo}

\subsection{Resultados}

\begin{table}[H]
  \centering
  \begin{tabular}{@{}lr@{}}
    \toprule
    \textbf{Métrica} & \textbf{Promedio} \\
    \midrule
    \filas{tablas/metricas_global.tex}
    \bottomrule
  \end{tabular}
  \caption{Métricas promedio sobre las consultas con juicios.}\label{tab:global}
\end{table}

\begin{figure}[H]
  \centering
  \begin{tikzpicture}
    \begin{axis}[xlabel={Posición en el ranking}, ylabel={Precisión promedio}, ymin=0, ymax=1,
      xtick={1,...,10}]
      \addplot[blue!70!black, mark=*] table[x=posicion, y=precision, col sep=comma]{datos/curva_precision.csv};
    \end{axis}
  \end{tikzpicture}
  \caption{Precisión promedio acumulada por posición.}
\end{figure}

\begin{figure}[H]
  \centering
  \begin{tikzpicture}
    \begin{axis}[ybar, bar width=5pt, ymin=0, ymax=1, legend pos=outer north east,
      xtick=data, xticklabels from table={datos/por_consulta.csv}{query_id},
      x tick label style={font=\scriptsize}, ylabel={Valor}]
      \addplot[fill=blue!60]   table[x=n, y=p,    col sep=comma]{datos/por_consulta.csv};
      \addplot[fill=orange!80] table[x=n, y=ap,   col sep=comma]{datos/por_consulta.csv};
      \addplot[fill=teal!70]   table[x=n, y=ndcg, col sep=comma]{datos/por_consulta.csv};
      \legend{P@10, AP, nDCG@10}
    \end{axis}
  \end{tikzpicture}
  \caption{Métricas por consulta.}
\end{figure}

\begin{table}[H]
  \centering
  \begin{tabular}{@{}ll@{}}
    \toprule
    \textbf{Id} & \textbf{Consulta} \\
    \midrule
    \filas{tablas/consultas.tex}
    \bottomrule
  \end{tabular}
  \caption{Consultas evaluadas.}
\end{table}

\begin{table}[H]
  \centering
  \begin{tabular}{@{}lrrrrr@{}}
    \toprule
    \textbf{Id} & \textbf{P@10} & \textbf{Recall@10} & \textbf{AP} & \textbf{nDCG@10} & \textbf{RR} \\
    \midrule
    \filas{tablas/metricas_consulta.tex}
    \bottomrule
  \end{tabular}
  \caption{Detalle por consulta (``--'' indica métrica no definida).}
\end{table}

\subsection{Limitaciones}

Los documentos sin juzgar se tratan como no relevantes, y las medidas habituales dejan de
ser robustas cuando los juicios son muy incompletos~\cite{buckley2004}; por eso se reporta
la cobertura de juicios. Además, Recall@$k$ se
calcula sobre el conjunto de documentos juzgados; si solo se juzga el top-$k$ vale 1 por
construcción y no debe presentarse como resultado. \pendiente{indicar la cobertura real de
juicios y cuántos evaluadores participaron por consulta}

% =========================================================
\section{Conclusiones}
% =========================================================

\subsection{¿Se cumplió la necesidad de información inicial?}
\pendiente{responder con las métricas de la Tabla~\ref{tab:global}, citando la necesidad de
información definida en Progra00 y qué documentos concretos la satisfacen o no}

\subsection{Conclusiones del grupo}
Mínimo 20; cada integrante aporta cerca de 7 desde su parte y Persona C las integra.
Cada conclusión debe apoyarse en un número, un gráfico o una observación de uso.

\begin{enumerate}[label=\textbf{C\arabic*.}, leftmargin=*]
  % Persona A: búsqueda y resultados (7)
  \item \pendiente{A-1: efecto de la búsqueda avanzada sobre los resultados}
  \item \pendiente{A-2}
  \item \pendiente{A-3}
  \item \pendiente{A-4}
  \item \pendiente{A-5}
  \item \pendiente{A-6}
  \item \pendiente{A-7}
  % Persona B: evaluación (7)
  \item \pendiente{B-1: dificultad de conseguir evaluadores y qué se aprendió}
  \item \pendiente{B-2}
  \item \pendiente{B-3}
  \item \pendiente{B-4}
  \item \pendiente{B-5}
  \item \pendiente{B-6}
  \item \pendiente{B-7}
  % Persona C: métricas (7)
  \item \pendiente{C-1: comparar P@10 con MAP y qué dice la diferencia}
  \item \pendiente{C-2: consultas con mejor y peor desempeño y por qué}
  \item \pendiente{C-3: qué aporta nDCG frente a las métricas binarias}
  \item \pendiente{C-4: MRR y consultas donde basta un buen resultado}
  \item \pendiente{C-5: acuerdo entre evaluadores ($\kappa$) y su efecto en las métricas}
  \item \pendiente{C-6: sesgo de tratar lo no juzgado como no relevante}
  \item \pendiente{C-7: desarrollo adicional y su efecto medido}
\end{enumerate}

% =========================================================
\begin{thebibliography}{9}
% Referencias verificadas el 24/11/2026 contra DOI, NIST, ACM y editoriales.
\bibitem{manning2008} C.~D. Manning, P.~Raghavan y H.~Schütze, \emph{Introduction to Information Retrieval}, cap.~8 (pp.~139--161), Cambridge University Press, 2008.
\bibitem{jarvelin2002} K.~Järvelin y J.~Kekäläinen, ``Cumulated gain-based evaluation of IR techniques'', \emph{ACM Transactions on Information Systems}, 20(4), pp.~422--446, 2002.
\bibitem{voorhees1999} E.~M. Voorhees, ``The TREC-8 Question Answering Track Report'', en \emph{Proc. TREC-8 (1999)}, NIST SP 500-246, pp.~77--82, 2000.
\bibitem{cohen1960} J.~Cohen, ``A coefficient of agreement for nominal scales'', \emph{Educational and Psychological Measurement}, 20(1), pp.~37--46, 1960.
\bibitem{landis1977} J.~R. Landis y G.~G. Koch, ``The measurement of observer agreement for categorical data'', \emph{Biometrics}, 33(1), pp.~159--174, 1977.
\bibitem{sanderson2010} M.~Sanderson, ``Test collection based evaluation of information retrieval systems'', \emph{Foundations and Trends in Information Retrieval}, 4(4), pp.~247--375, 2010.
\bibitem{buckley2004} C.~Buckley y E.~M. Voorhees, ``Retrieval evaluation with incomplete information'', en \emph{Proc. SIGIR}, pp.~25--32, 2004.
\bibitem{voorhees2000} E.~M. Voorhees, ``Variations in relevance judgments and the measurement of retrieval effectiveness'', \emph{Information Processing \& Management}, 36(5), pp.~697--716, 2000.
\end{thebibliography}

\end{document}
